\begin{landscape}
\begingroup
\fontsize{10}{13}\selectfont
\fancyhead[L]{\small POSTHUMOUS PAPERS AND CORRESPONDENCE}
\renewcommand{\headrulewidth}{0pt}
\par\medskip
\pdfbookmark[2]{Posthumous Papers VII [2]}{pallas-VII-fragment2}
\begin{center}[2.]\end{center}
\makeatletter
\providecommand{\VIItwoHistoricalNote}[1]{\begingroup
\renewcommand{\thefootnote}{*)}
\long\def\@makefntext##1{\noindent [*)\ ##1}
\footnote{#1}\endgroup}
\makeatother
\par\noindent\hspace*{\parindent}[To facilitate the calculation of the amounts of perturbation for the individual oppositions, tables can be
constructed as follows:
\par\noindent\hspace*{\parindent}First combine the perturbation equations with the arguments
\[
\begin{gathered}
2\text{\mdseries\AstroSymbols ♃}-\text{\mdseries\AstroSymbols ⚴},\quad
\text{\mdseries\AstroSymbols ♃},\quad\text{\mdseries\AstroSymbols ⚴},\quad
\text{\mdseries\AstroSymbols ♃}-2\text{\mdseries\AstroSymbols ⚴},\quad
2\text{\mdseries\AstroSymbols ♃}-3\text{\mdseries\AstroSymbols ⚴}\ \text{etc.}\\
3\text{\mdseries\AstroSymbols ♃}-2\text{\mdseries\AstroSymbols ⚴},\quad
4\text{\mdseries\AstroSymbols ♃}-3\text{\mdseries\AstroSymbols ⚴}\ \text{etc.,}
\end{gathered}
\]
\noindent into one group. Setting
\[
2\text{\mdseries\AstroSymbols ♃}-\text{\mdseries\AstroSymbols ⚴}=A_1,\qquad
\text{\mdseries\AstroSymbols ⚴}-\text{\mdseries\AstroSymbols ♃}=x
\]
\noindent puts them in the form
\[
A_1+nx+G^{(n)},
\]
\noindent where \(n\) denotes the positive and negative integers, including zero, and \(G^{(n)}\) the corresponding constant.
We may write:
\begingroup\small
\[
\begin{aligned}
&g\sin(A_1+G)+g'\sin(A_1+x+G')+g''\sin(A_1+2x+G'')+\cdots\\
&\qquad+h'\sin(A_1-x+H')+h''\sin(A_1-2x+H'')+\cdots\\
{}={}&\sin A_1\,\mathord{.}\,(a+a'\cos x+a''\cos2x+\cdots+b'\sin x+b''\sin2x+\cdots)\\
&+\cos A_1\,\mathord{.}\,(c+c'\cos x+c''\cos2x+\cdots+d'\sin x+d''\sin2x+\cdots),
\end{aligned}
\]
\endgroup
\noindent where
\begingroup\scriptsize
\[
\begin{array}{l l l}
a=g\cos G,&a'=g'\cos G'+h'\cos H',&a''=g''\cos G''+h''\cos H'',\quad\text{etc.}\\
&b'=-g'\sin G'+h'\sin H',&b''=-g''\sin G''+h''\sin H'',\quad\text{etc.}\\
c=g\sin G,&c'=g'\sin G'+h'\sin H',&c''=g''\sin G''+h''\sin H'',\quad\text{etc.}\\
&d'=g'\cos G'-h'\cos H',&d''=g''\cos G''-h''\cos H'',\quad\text{etc.}
\end{array}
\]
\endgroup
\par\noindent\hspace*{\parindent}Now if we determine \(r\) and \(\varphi\) from the equations
\[
\begin{aligned}
\left.\begin{aligned}
a+a'\cos x+a''\cos2x+\cdots\\
{}+b'\sin x+b''\sin2x+\cdots
\end{aligned}\right\}&=P=r\cos\varphi\\[1ex]
\left.\begin{aligned}
c+c'\cos x+c''\cos2x+\cdots\\
{}+d'\sin x+d''\sin2x+\cdots
\end{aligned}\right\}&=Q=r\sin\varphi,
\end{aligned}
\]
\SourceLeaf{574}
\noindent we obtain for the whole group the expression
\begingroup\small
\[
\left.\begin{aligned}
g\sin(A_1+G)+g'\sin(A_1+x+G')+g''\sin(A_1+2x+G'')+\cdots\\
{}+h'\sin(A_1-x+H')+h''\sin(A_1-2x+H'')+\cdots
\end{aligned}\right\}=r\sin(A_1+\varphi).
\]
\endgroup
\par\noindent\hspace*{\parindent}The \textls[100]{first operation}\VIItwoHistoricalNote{See the following letter.]} now consists in calculating \(a^{(n)}\), \(b^{(n)}\), \(c^{(n)}\), \(d^{(n)}\) from \(g^{(n)}\), \(h^{(n)}\),
\(G^{(n)}\), \(H^{(n)}\). The \textls[100]{second} consists in calculating \(P\) and \(Q\) for the values \(x=0^0\), \(15^0\), \(30^0\), etc., and the
\textls[100]{third} in calculating \(r\) and \(\varphi\).
\par\noindent\hspace*{\parindent}The remaining perturbation equations are to be combined into groups in the same way by taking
as a fixed definition
\[
x=\text{\mdseries\AstroSymbols ⚴}-\text{\mdseries\AstroSymbols ♃},
\]
\noindent but using
\begingroup\footnotesize
\[
\begin{array}{l|l|l|l}
A_2=3\text{\mdseries\AstroSymbols ♃}-\text{\mdseries\AstroSymbols ⚴}&
A_5=8\text{\mdseries\AstroSymbols ♃}-3\text{\mdseries\AstroSymbols ⚴}&
A_8=13\text{\mdseries\AstroSymbols ♃}-5\text{\mdseries\AstroSymbols ⚴}&
A_{11}=17\text{\mdseries\AstroSymbols ♃}-6\text{\mdseries\AstroSymbols ⚴}\\
A_3=5\text{\mdseries\AstroSymbols ♃}-2\text{\mdseries\AstroSymbols ⚴}&
A_6=10\text{\mdseries\AstroSymbols ♃}-4\text{\mdseries\AstroSymbols ⚴}&
A_9=15\text{\mdseries\AstroSymbols ♃}-6\text{\mdseries\AstroSymbols ⚴}&
A_{12}=20\text{\mdseries\AstroSymbols ♃}-8\text{\mdseries\AstroSymbols ⚴}\\
A_4=6\text{\mdseries\AstroSymbols ♃}-2\text{\mdseries\AstroSymbols ⚴}&
A_7=11\text{\mdseries\AstroSymbols ♃}-4\text{\mdseries\AstroSymbols ⚴}&
A_{10}=16\text{\mdseries\AstroSymbols ♃}-6\text{\mdseries\AstroSymbols ⚴}&
A_{13}=21\text{\mdseries\AstroSymbols ♃}-8\text{\mdseries\AstroSymbols ⚴}
\end{array}
\]
\endgroup
\noindent for the other definitions.
\par\noindent\hspace*{\parindent}In this way the table printed as Note [4.] was constructed. For each \(A\) it gives the quantities
\(r\) and \(\varphi\) for the argument \(x\) at intervals of \(15^0\).
\par\noindent\hspace*{\parindent}The group of perturbation equations with the arguments \(\text{\mdseries\AstroSymbols ♃}-\text{\mdseries\AstroSymbols ⚴}\), \(2\text{\mdseries\AstroSymbols ♃}-2\text{\mdseries\AstroSymbols ⚴}\), \(3\text{\mdseries\AstroSymbols ♃}-3\text{\mdseries\AstroSymbols ⚴}\) etc. can be calculated directly
for the argument \(x\). At the beginning of the table under the heading \(A_0=0\), their values are
given.]
\clearpage
\endgroup
\end{landscape}
